% 第2章 2.5.2：伴随与像、核的正交分解关系图（示意）。
% H_1 = ker A ⊕ cl(Ran A^*)，H_2 = ker A^* ⊕ cl(Ran A)。
% A 把 ker A 送到 0，像落在 Ran A 中；A^* 把 ker A^* 送到 0，像落在 Ran A^* 中。
\begin{tikzpicture}[x=1mm,y=1mm,font=\figurefont\footnotesize,text=FigureInk,
  blk/.style={draw=FigureBlue,line width=.8pt,rounded corners=1.8pt,minimum width=34mm,align=center,inner sep=1.5mm},
  note/.style={font=\FigureNoteFont\fontsize{7.5}{9}\selectfont,text=FigureMuted},
  arr/.style={line width=.8pt,-{Stealth[length=1.8mm]}}]
\node[blk,fill=FigureBlueFill,minimum height=15mm] (RAs) at (0,10) {$\overline{\operatorname{Ran}A^*}$\\[-.3mm]{\FigureNoteFont\fontsize{7.5}{9}\selectfont 可由输出测试拉回}};
\node[blk,fill=white,minimum height=11mm] (KA) at (0,-6) {$\ker A$\\[-.3mm]{\FigureNoteFont\fontsize{7.5}{9}\selectfont 输出中看不见}};
\node[blk,fill=FigureBlueFill,minimum height=15mm] (RA) at (62,10) {$\overline{\operatorname{Ran}A}$\\[-.3mm]{\FigureNoteFont\fontsize{7.5}{9}\selectfont 可被输出逼近}};
\node[blk,fill=white,minimum height=11mm] (KAs) at (62,-6) {$\ker A^*$\\[-.3mm]{\FigureNoteFont\fontsize{7.5}{9}\selectfont 输出永远到不了}};
\node at (0,1.6) {$\perp$}; \node at (62,1.6) {$\perp$};
\node[above=1mm] at (RAs.north) {$H_1$}; \node[above=1mm] at (RA.north) {$H_2$};
\draw[arr,FigureBlue] ([yshift=3mm]RAs.east)--node[above] {$A$} ([yshift=3mm]RA.west);
\draw[arr,FigureGreen] ([yshift=-3mm]RA.west)--node[below] {$A^*$} ([yshift=-3mm]RAs.west -| RAs.east);
\node[circle,draw=FigureStroke,line width=.6pt,inner sep=.8mm] (Z) at (31,-6) {$0$};
\draw[arr,FigureBlue] (KA.east)--node[above,font=\figurefont\scriptsize] {$A$} (Z.west);
\draw[arr,FigureGreen] (KAs.west)--node[above,font=\figurefont\scriptsize] {$A^*$} (Z.east);
\end{tikzpicture}
